\documentclass[10pt,a4paper]{amsart}
\usepackage[margin=19mm]{geometry}
\usepackage{amsmath,amssymb,mathtools}
\usepackage[scr=rsfs,cal=euler]{mathalfa}
\usepackage{xfrac}
\usepackage[unicode,hidelinks,bookmarks=false]{hyperref}
\newcommand{\ie}{i.e.\ }
\newcommand{\cf}{cf.\ }
\newcommand{\resp}{respectively\ }
\newcommand{\Subsection}{\subsection}
\providecommand{\nobreakdash}{\nobreak\hbox{-}}
\setlength{\emergencystretch}{2em}
\multlinegap0pt
\usepackage{mathtools}
\usepackage{amssymb}
\usepackage{bm}
\newtheorem{theorem}{Theorem}
\def \N{\mathbb N}
\def \Z{\mathbb Z}
\def \k{\bm{k}}
\title{{\bf Almost global large deviations principle for the KdV equation}}
\author{Riccardo Berforini D'Aquino, Ricardo Grande}
\date{Source: arXiv:2605.00484v1 (CC BY 4.0)}
\begin{document}
\maketitle

\noindent\emph{Source and licence.} {\bf Almost global large deviations principle for the KdV equation}. Version arXiv:2605.00484v1 (CC BY 4.0), distributed by arXiv under the Creative Commons Attribution 4.0 International licence, http://creativecommons.org/licenses/by/4.0/, as recorded in the arXiv licence metadata for this exact version and shown on the abstract page https://arxiv.org/abs/2605.00484v1. Full licence notice: \texttt{licenses/arxiv-2605.00484-CC-BY-4.0.txt}.\par\smallskip

\noindent\textit{Editorial scope.} Counted unit: the article's theorem 2 (source0.tex lines 3459--3484). The article's complete original proof of it is delivered with it (source0.tex lines 3485--3634, one \texttt{proof} environment of 150 lines). No mathematics is written, repaired or re-derived: the statement and the proof are the article's own lines, in the article's own notation, and no retained line is substituted or re-flowed.  No further source-local context is needed: every cross-reference of every retained line resolves inside the two delivered spans.  Nothing was omitted from any retained span, so the package carries no omission record.  No retained line contains a citation, so the wrapper carries no bibliography and no bibliography entry.  No retained line contains a figure environment, an image-inclusion command or drawing code, so no image is delivered and the package declares no asset dependency.  Editorial formatting only, in addition: an AMS \texttt{amsart} preamble that restates the article's own declarations and macros used by the retained lines, namely the environments \texttt{theorem} on separate counters, together with the standard packages those lines need.  The retained environments print in this document as Theorem 1 in order of appearance, whereas the article prints the same items with the numbering of its own full document.  The complete native source of the article as distributed in the arXiv e-print for this exact version is bundled unchanged as \texttt{sources/199567/source0.tex}.
\begin{theorem} \label{thm: system for resonant 2}
Given $n\in\N$, consider the system of $n$ equations
\begin{equation}
\begin{cases}
\alpha_1 k_1 + \alpha_2 k_2+\ldots+ \alpha_n k_n & = 0 \\
\alpha_1 k_1^3 + \alpha_2 k_2^3 + \ldots+ \alpha_n k_n^3 & = 0 \\
& \vdots \\
\alpha_1 k_1^{2n-1} + \alpha_2 k_2^{2n-1} +\ldots + \alpha_n k_n^{2n-1} & =0,
\end{cases}
\label{system 3}
\end{equation}
where $\alpha_j \in \N$, $j=1,\ldots, n$. We require that $k_j \neq 0$ for all $j=1,\ldots,n$. Then:

\begin{itemize}
\item If $S = \alpha_1 + \alpha_2 + \ldots + \alpha_n$ is odd, there is no integer solution to \eqref{system 3}. 

\item If $S$ is even, all the solutions $(k_1,k_2,\ldots,k_n)$ are \emph{weightedly paired}, meaning that if we build the vector 
\[
(k_1',\ldots,k_S')=(k_{1,1},\ldots,k_{1,\alpha_1},k_{2,1},\ldots,k_{2,\alpha_2},\ldots,k_{n,1},\ldots,k_{n,\alpha_n}) \in \Z^{S},
\]
where $k_{j,1}=k_{j,2}=\ldots=k_{j,\alpha_j}=k_j$, there exists a permutation $\sigma \in Sym(S)$ such that 
\[
k_{\sigma(2i-1)}'+k_{\sigma(2i)}'=0 \qquad \text{for $i=1,\ldots,\frac{S}{2}.$}
\]
\end{itemize}
\end{theorem}

\begin{proof}
First, we write \eqref{system 3} as a system:
\[
\begin{bmatrix}
k_1 & k_2  & \dots  & k_n \\
k_1^3 & k_2^3 & \dots  & k_n^3 \\
\vdots & \vdots & \ddots & \vdots \\
k_1^{2n-1} & k_2^{2n-1} & \dots  & k_n^{2n-1}
\end{bmatrix} 
\cdot
\begin{bmatrix}
\alpha_1 \\
\alpha_2 \\
\vdots \\
\alpha_n
\end{bmatrix}
=
\begin{bmatrix}
0 \\
0 \\
\vdots \\
0
\end{bmatrix}.
\]
Since $\alpha_1,\ldots,\alpha_n\neq 0$, we must have a singular matrix: 
\begin{equation} 
\begin{split}
0 & = \mathrm{det}
\begin{bmatrix}
k_1 & k_2  & \dots  & k_n \\
k_1^3 & k_2^3 & \dots  & k_n^3 \\
\vdots & \vdots & \ddots & \vdots \\
k_1^{2n-1} & k_2^{2n-1} & \dots  & k_n^{2n-1}
\end{bmatrix}
= k_1k_2\dots k_n \cdot \mathrm{det}
\begin{bmatrix}
1 & 1  & \dots  & 1 \\
k_1^2 & k_2^2 & \dots  & k_n^2 \\
\vdots & \vdots & \ddots & \vdots \\
k_1^{2n-2} & k_2^{2n-2} & \dots  & k_n^{2n-2}
\end{bmatrix} \\
& = k_1k_2\dots k_n \prod_{1 \le i < j \le n} (k_i^2-k_j^2) = k_1k_2\dots k_n \prod_{1 \le i < j \le n}(k_i-k_j)(k_i+k_j)=0,
\label{vandermonde2}
\end{split}
\end{equation}
where we used the explicit formula for the determinant of the Vandermonde matrix.
This implies that there exist $i,j$ such that $k_i+k_j=0$ or $k_i-k_j=0$. 
%In the former case we say that $(i,j)$ is a \textit{good pair}, in the latter one we say that it is a \textit{bad pair}. 
The rest of the proof proceeds by induction, in particular we will prove that if the theorem is true for a system of dimension $n-2$ and $n-1$, then it is true for $n$. We need two base cases.

\smallskip

\noindent {\sc Base Case 1}: If $n=1$, then $\alpha_1 k_1 \neq 0$ by hypothesis but \eqref{system 3} forces $\alpha_1 k_1 = 0$, hence there is no solution.

\smallskip

\noindent {\sc Base Case 2:} If $n=2$, \eqref{system 3}  becomes:
\begin{equation}
\begin{cases}
\alpha_1 k_1 + \alpha_2 k_2 = 0 \\
\alpha_1 k_1^3 + \alpha_2 k_2^3 = 0, 
\end{cases} \qquad \qquad \mbox{with}\ \alpha_1,\alpha_2 \in \N.
\label{finale 1,1}
\end{equation}
If we have a solution $(k_1,k_2)$ to \eqref{finale 1,1} with $k_1k_2 \neq 0$, then \eqref{vandermonde2} implies $k_1^2-k_2^2=0$. However, we cannot have $k_1=k_2$, since otherwise \eqref{finale 1,1} implies that $\alpha_1 + \alpha_2 = 0$, which is not possible. Therefore $k_1=-k_2 \neq 0$. Now we distinguish the two cases presented in the statement:
\begin{itemize}
    \item If $S=\alpha_1+\alpha_2$ is odd, there exist no solutions by Base Case 1. 
    \item If $S=\alpha_1+\alpha_2$ is even, then $(k_1,k_2)$ is a solution if and only if $\alpha_1=\alpha_2,$ which means that $(k_1,k_2)$ is weightedly paired.
\end{itemize}
% \textbf{Base case 2}:
% \begin{equation}
% \begin{cases}
% \alpha_1 k_1 + \alpha_2 k_2 + \alpha_3 k_3=0 \\
% \alpha_1 k_1^3 + \alpha_2 k_2^3 + \alpha_3 k_3^3 = 0 \\
% \alpha_1 k_1^5 + \alpha_2 k_2^5 + \alpha_3 k_3^5 = 0,
% \end{cases}
% \label{finale 2,1}
% \end{equation}
% with $\alpha_1,\alpha_2,\alpha_3 >0$. Here there are three possibilities: 
% \begin{itemize}
% \item If $k_i-k_j=0$, up to a permutation of indices we can assume $k_2=k_3$ and we arrive to the base case $1$
% \[
% \begin{cases}
% \alpha_1' k_1 + \alpha_2' k_2 = 0 \\
% \alpha_1' k_1^3 + \alpha_2' k_2^3 = 0,
% \end{cases}
% \]
% with $\alpha_2' = \alpha_2+\alpha_3$ and $\alpha_1' = \alpha_1$. But we know that this has a solution only if $\alpha_1'=\alpha_2'$ and the solution is $k_1=-k_2$. Then the solution of \eqref{finale 2,1} is weightedly paired. 
% Notice that in this case, if $S$ were odd, there would be no solution since 
% \[
% S = \alpha_1 + \alpha_2 + \alpha_3 = \alpha_1' + \alpha_2' = S'
% \]
% and the only case with solutions is when $S = \alpha_1' + \alpha_2' = 2 \alpha_1'$ which is even.

% \item If $k_i+k_j=0$ with $\alpha_i-\alpha_j \neq 0$, up to a permutation of indices we arrive to the base case $1$
% \[
% \begin{cases}
% \alpha_1' k_1 + \alpha_2' k_2 = 0 \\
% \alpha_1' k_1^3 + \alpha_2' k_2^3 = 0,
% \end{cases}
% \]
% with $\alpha_2' = \alpha_2-\alpha_3 > 0$ and $\alpha_1' = \alpha_1$. But we know that this has a solution only if $\alpha_1'=\alpha_2'$ and the solution is $k_1=-k_2$. Then the solution of \eqref{finale 2,1} is weightedly paired. Notice that, setting $S' = \alpha_1' + \alpha_2'$, we have $S = S' + 2 \alpha_3$ and so $S-S' \equiv 0 \quad (mod \, \, 2).$

% \item If $k_i+k_j=0$ with $\alpha_i-\alpha_j = 0$, up to a permutation of indices we arrive to $k_1=0$ which is not possible.
% \end{itemize}

\medskip

\noindent {\sc General Case:} Next we consider \eqref{system 3} with $n\geq 3$.
Let $\k = (k_1,\ldots,k_n)$ be a solution to \eqref{system 3} with $\prod_{j=1}^n k_j \neq 0$.
By \eqref{vandermonde2}, there are three possibilities:

\begin{itemize}
\item If $k_i-k_j=0$, up to a permutation of indices we can assume $k_{n-1}=k_n$ and we arrive to 
\begin{equation}\label{system_ind}
\begin{cases}
\alpha_1' k_1 + \alpha_2' k_2+\ldots+ \alpha_{n-1}' k_{n-1} & = 0 \\
\alpha_1' k_1^3 + \alpha_2' k_2^3 + \ldots+ \alpha_{n-1}' k_{n-1}^3 & = 0 \\
& \vdots \\
\alpha_1' k_1^{2n-3} + \alpha_2' k_2^{2n-3} +\ldots + \alpha_{n-1}' k_{n-1}^{2n-3} & =0,
\end{cases}
\end{equation}
with $\alpha_{n-1}' = \alpha_{n-1}+\alpha_{n}$ and $\alpha_i' = \alpha_i$ for $i=1,\ldots,n-2$. 
By the inductive hypothesis, since $S'=S$, if $S$ is even, the system \eqref{system_ind} has weightedly paired solutions, therefore $\k$ is weightedly paired. If $S$ is odd we have no solutions to \eqref{system_ind}, and thus $\k$ doesn't solve \eqref{system 3}. 

\item If $k_i+k_j=0$ with $\alpha_i-\alpha_j \neq 0$, up to a permutation of indices we arrive to 
\begin{equation} \label{system_ind2}
\begin{cases}
\alpha_1' k_1 + \alpha_2' k_2+\ldots+ \alpha_{n-1}' k_{n-1} & = 0 \\
\alpha_1' k_1^3 + \alpha_2' k_2^3 + \ldots+ \alpha_{n-1}' k_{n-1}^3 & = 0 \\
& \vdots \\
\alpha_1' k_1^{2n-3} + \alpha_2' k_2^{2n-3} +\ldots + \alpha_{n-1}' k_{n-1}^{2n-3} & =0,
\end{cases}
\end{equation}
with $\alpha_{n-1}' = \alpha_{n-1}-\alpha_{n} \in \N$ (without loss of generality) and $\alpha_i' = \alpha_i$ for $i=1,\ldots,n-2$. Notice that the indices we have removed are already weightedly paired. 
Since $S'=S-2\alpha_n$, $S-S' \equiv 0 \, (\mathrm{mod} \, \, 2)$, hence by the inductive hypothesis the system \eqref{system_ind2} has weightedly paired solutions if $S$ is even (and the same holds for $\k$), and no solutions if $S$ is odd. 

\item If $k_i+k_j=0$ with $\alpha_i-\alpha_j = 0$, up to a permutation of indices can assume $k_{n-1}+k_n=0$ and we arrive to 
\begin{equation} \label{syst_ind3}
\begin{cases}
\alpha_1' k_1 + \alpha_2' k_2+\ldots+ \alpha_{n-2}' k_{n-2} & = 0 \\
\alpha_1' k_1^3 + \alpha_2' k_2^3 + \ldots+ \alpha_{n-2}' k_{n-2}^3 & = 0 \\
& \vdots \\
\alpha_1' k_1^{2n-5} + \alpha_2' k_2^{2n-5} +\ldots + \alpha_{n-2}' k_{n-2}^{2n-5} & =0.
\end{cases}
\end{equation}
with $\alpha_j'=\alpha_j$ for $j=1,\ldots,n-2$.
Notice that $k_i$ and $k_j$ that we have removed are already weightedly paired. Again $S-S' \equiv 0 \quad (mod \, \, 2)$, hence by the inductive hypothesis the system \eqref{syst_ind3} has weightedly paired solutions if $S$ is even (and the same holds for $\k$), and no solutions if $S$ is odd. 
\end{itemize}
\end{proof}

\end{document}
